% Generated draft: requires mathematical and source review.
\documentclass[11pt]{article}
\usepackage[margin=25mm]{geometry}
\usepackage{fontspec}
\setmainfont{Latin Modern Roman}
\usepackage{amsmath,amssymb,mathtools,unicode-math}
\setmathfont{Latin Modern Math}
\usepackage{enumitem,needspace,xurl,hyperref}
\setlength{\parindent}{0pt}
\setlength{\parskip}{5pt}
\setlength{\emergencystretch}{3em}
\newcommand{\N}{\mathbb{N}}
\newcommand{\Z}{\mathbb{Z}}
\newcommand{\Q}{\mathbb{Q}}
\newcommand{\R}{\mathbb{R}}
\newcommand{\C}{\mathbb{C}}
\newcommand{\F}{\mathbb{F}}
\newcommand{\A}{\mathbb{A}}
\newcommand{\PP}{\mathbb{P}}
\newcommand{\E}{\mathbb{E}}
\newcommand{\eps}{\varepsilon}
\newcommand{\dd}{\,\mathrm{d}}
\begin{document}
\section*{UnsolvedMath ID: 701644}
\textbf{Source problem code:} catalogue entry
\subsection{Definitions and mathematical statement}
Let \(r\geq 2\) be a finite integer, let \(F_r\) be the free group on \(r\) generators, and write \(F_r'=[F_r,F_r]\) and \(F_r''=[F_r',F_r']\). Let \(M\cong F_r/F_r''\), so that \(M\) is the free metabelian group of rank \(r\). A metabelian group is a group whose commutator subgroup is abelian.

An ordered tuple \((b_1,\ldots,b_r)\in M^r\) is a basis of \(M\) if the homomorphism \(F_r/F_r''\to M\) sending the standard generators to \(b_1,\ldots,b_r\) is an isomorphism. Define
\[
\mathcal B_r(M)=\{(b_1,\ldots,b_r)\in M^r:(b_1,\ldots,b_r)\text{ is a basis of }M\}.
\]
An element is primitive if it belongs to a basis. Thus define
\[
P(M)=\{g\in M:\exists b_2,\ldots,b_r\in M\text{ with }(g,b_2,\ldots,b_r)\in\mathcal B_r(M)\}.
\]
The existential quantifier in this displayed definition is an ordinary mathematical quantifier; it does not assume that \(\mathcal B_r(M)\) is first-order definable.

\textbf{Questions (E. I. Timoshenko, Kourovka Notebook Problem 19.96).}
\begin{enumerate}
\item[(a)] Is \(P(M)\) a first-order definable subset of \(M\)?
\item[(b)] Is \(\mathcal B_r(M)\) a first-order definable subset of \(M^r\)?
\end{enumerate}
Here the standard group language is \(\mathcal L_{\mathrm{grp}}=\{\cdot,{}^{-1},1\}\), with variables ranging over elements of \(M\). Under this language convention, definability in (a) means the existence of a formula \(\varphi(x;\mathbf z)\) and an admissible finite parameter tuple \(\mathbf a\) from \(M\) such that
\[
\forall g\in M,\qquad g\in P(M)\ \Longleftrightarrow\ M\models\varphi(g;\mathbf a).
\]
In (b), it means the existence of a formula \(\psi(x_1,\ldots,x_r;\mathbf w)\) and an admissible finite parameter tuple \(\mathbf c\) from \(M\) such that
\[
\forall(b_1,\ldots,b_r)\in M^r,\qquad
(b_1,\ldots,b_r)\in\mathcal B_r(M)
\ \Longleftrightarrow\ M\models\psi(b_1,\ldots,b_r;\mathbf c).
\]
\textbf{Unresolved qualifications.} The supplied statement does not explicitly specify the first-order language or whether parameters are allowed. The formulas above explain the standard group-language interpretation without resolving the parameter convention: for parameter-free definability the parameter tuples must be empty, whereas definability with parameters permits finite tuples from \(M\). Part (b) is represented by ordered basis tuples; the corresponding family is invariant under permutation of its coordinates. The source does not explicitly specify this encoding or impose a uniformity requirement as the rank varies.
\subsection{Short English statement}
For a free metabelian group of finite rank \(r\geq 2\), E. I. Timoshenko asks whether its primitive elements and its bases can each be described by a first-order formula. Bases are represented here by ordered \(r\)-tuples. The source leaves the language, parameter convention, and any uniformity across ranks unspecified.
\subsection{Sources}
\begin{itemize}
\item Local UnsolvedMath registry, permanent ID \(701644\), exact source problem code \texttt{catalogue entry}, title \emph{Kourovka Notebook Problem 19.96}. The supplied record attributes the problem to E. I. Timoshenko, with proposed year 2018.
\item Supplied source attribution: \emph{The Kourovka Notebook (Unsolved Problems in Group Theory)}, 19th issue (2018), Problem 19.96; also identified in the record with the 21st issue edition (2026), \texttt{arXiv:1401.0300v48}. The collection description names E. I. Khukhro and V. D. Mazurov as editors. These bibliographic details are supplied attribution and have not been independently verified here.
\item The supplied registry labels the problem \texttt{open} and gives an update timestamp of \texttt{2026-10-16T14:00:00Z}. This is source metadata, not an independently verified assessment of the mathematical status.
\end{itemize}
\end{document}
